% Part: computability
% Chapter: recursive-functions
% Section: halting-problem

\documentclass[../../../include/open-logic-section]{subfiles}

\begin{document}

\olfileid{cmp}{rec}{hlt}
\olsection{থামার সমস্যা}

সাধারণভাবে \emph{থামার সমস্যা} হলো এই সিদ্ধান্ত নেওয়ার সমস্যা: একটি
গণনসাধ্য অপেক্ষকের নির্দেশন~$e$ (যেমন কোনো প্রোগ্রাম) এবং একটি
সংখ্যা~$n$ দেওয়া থাকলে, নিবেশ~$n$-এ অপেক্ষকটির গণনা থামে, অর্থাৎ কোনো
ফল উৎপন্ন করে, কি না। অ্যালান টুরিংয়ের বিখ্যাত প্রমাণ অনুসারে এই সমস্যার
সমাধান নিজেই কোনো গণনসাধ্য অপেক্ষক দিয়ে করা যায় না; অর্থাৎ অপেক্ষক
\[
h(e, n) =
\begin{cases}
1 & \text{যদি নিবেশ $n$-এ গণনা $e$ থামে}\\
0 & \text{অন্যথায়,}
\end{cases}
\]
গণনসাধ্য নয়।

আংশিক পুনরাবৃত্ত অপেক্ষকের প্রসঙ্গে কোনো প্রোগ্রামের নির্দেশনের ভূমিকা
পালন করতে পারে ক্লিনির স্বাভাবিক রূপের উপপাদ্যে দেওয়া সূচক~$e$। $f$
কোনো আংশিক পুনরাবৃত্ত অপেক্ষক হলে, স্বাভাবিক রূপের উপপাদ্যের সমীকরণ যে
$e$-এর জন্য সত্য, সেটিই~$f$-এর একটি সূচক। কোনো
সংখ্যা~$e$ দেওয়া থাকলে স্বাভাবিক রূপের উপপাদ্য বলে যে
\[
\cfind{e}(x) \simeq U(\mu s \; T(e, x, s))
\]
আংশিক পুনরাবৃত্ত; এবং প্রত্যেক আংশিক পুনরাবৃত্ত $f\colon \Nat \to
\Nat$-এর জন্য এমন একটি $e \in \Nat$ আছে, যাতে সব~$x \in \Nat$-এর
ক্ষেত্রে $\cfind{e}(x) \simeq f(x)$। আসলে এমন প্রত্যেক $f$-এর জন্য শুধু
একটি নয়, অসীমসংখ্যক এমন~$e$ আছে। \emph{থামা-অপেক্ষক}~$h$-কে এভাবে
সংজ্ঞায়িত করা হয়:
\[
h(e, x) =
\begin{cases}
1 & \text{যদি $\cfind{e}(x) \fdefined$}\\
0 & \text{অন্যথায়।}
\end{cases}
\]
লক্ষ করুন, $\cfind{e}(x) \fundefined$ হলে $h(e, x) = 0$, আবার~$e$
আদৌ কোনো আংশিক পুনরাবৃত্ত অপেক্ষকের সূচক না হলেও তাই।

\begin{thm}
\ollabel{thm:halting-problem}
থামা-অপেক্ষক~$h$ আংশিক পুনরাবৃত্ত নয়।
\end{thm}

\begin{proof}
$h$ আংশিক পুনরাবৃত্ত হলে আমরা নিচের অপেক্ষকটি সংজ্ঞায়িত করতে পারতাম:
\[
d(y) =
\begin{cases}
1 & \text{যদি $h(y, y) = 0$}\\
\umin{x}{x \neq x} & \text{অন্যথায়।}
\end{cases}
\]
কোনো সংখ্যা $x$, $x \neq x$ পূরণ করে না বলে $\umin{x}{x \neq x}$-এর
অস্তিত্ব নেই; তাই $d(y) \fundefined$ যদি এবং কেবল যদি $h(y,y) \neq
0$। এই সংজ্ঞা থেকে পাওয়া যায়:
\begin{enumerate}
\item $d(y) \fdefined$ যদি এবং কেবল যদি $\cfind{y}(y) \fundefined$,
  অথবা $y$ কোনো আংশিক পুনরাবৃত্ত অপেক্ষকের সূচক নয়।
\item $d(y) \fundefined$ যদি এবং কেবল যদি $\cfind{y}(y) \fdefined$।
\end{enumerate}
$h$ আংশিক পুনরাবৃত্ত হলে $d$-ও আংশিক পুনরাবৃত্ত হতো। তাহলে ক্লিনির
স্বাভাবিক রূপের উপপাদ্য অনুসারে এর একটি সূচক~$e_d$ থাকত। $h(e_d, e_d)$-এর
মান বিবেচনা করুন। দুটি সম্ভাব্য ক্ষেত্র আছে—$0$ ও~$1$।
\begin{enumerate}
\item $h(e_d, e_d) = 1$ হলে $\cfind{e_d}(e_d) \fdefined$। কিন্তু
  $\cfind{e_d} \simeq d$, আর $d(e_d)$ সংজ্ঞায়িত যদি এবং কেবল যদি
  $h(e_d, e_d) = 0$। অতএব $h(e_d, e_d) \neq 1$।
\item $h(e_d, e_d) = 0$ হলে হয় $e_d$ কোনো আংশিক পুনরাবৃত্ত অপেক্ষকের
  সূচক নয়, অথবা সূচক হলেও $\cfind{e_d}(e_d) \fundefined$। কিন্তু আবার
  $\cfind{e_d} \simeq d$, আর $d(e_d)$ অসংজ্ঞায়িত যদি এবং কেবল যদি
  $\cfind{e_d}(e_d) \fdefined$।
\end{enumerate}
ফল হলো, শেষ পর্যন্ত $e_d$ কোনো আংশিক পুনরাবৃত্ত অপেক্ষকের সূচক হতে পারে
না। কিন্তু $h$ আংশিক পুনরাবৃত্ত হলে $d$-ও তাই হতো, এবং তখন তার সূচক
হিসেবে $e_d$-এর সংজ্ঞা গ্রহণযোগ্য হতো। তাই আমাদের সিদ্ধান্ত নিতে হবে,
$h$ আংশিক পুনরাবৃত্ত হতে পারে না।
\end{proof}

\end{document}
